\documentclass[11pt]{article}
\usepackage[margin=1in]{geometry}
\usepackage[T1]{fontenc}
\usepackage[utf8]{inputenc}
\usepackage{lmodern}
\usepackage{booktabs}
\usepackage{graphicx}
\usepackage{caption}
\usepackage{enumitem}
\usepackage{hyperref}
\usepackage{natbib}
\hypersetup{colorlinks=true, linkcolor=black, citecolor=black, urlcolor=blue}
\captionsetup{font=small, labelfont=bf}
\setlength{\tabcolsep}{4pt}
\graphicspath{{../../output/reviewed/llm_gpt-5-mini-2025-08-07/figures/}}
\newcommand{\webAdvice}{7.7}
\newcommand{\webNeverAdvise}{42}
\newcommand{\webMostlyAdvise}{3}
\newcommand{\webPairAgreement}{31}
\newcommand{\webEligiblePrompts}{60}
\newcommand{\poolLone}{0.0}
\newcommand{\poolLtwo}{29.4}
\newcommand{\poolLthree}{1.0}
\newcommand{\poolLfour}{27.9}
\newcommand{\poolLfive}{42.3}
\newcommand{\zemaRec}{20}
\newcommand{\zemaPoll}{2}
\newcommand{\lulaRec}{22}
\newcommand{\lulaPoll}{38}
\newcommand{\flavioRec}{23}
\newcommand{\flavioPoll}{31}
\newcommand{\caiadoRec}{14}
\newcommand{\caiadoPoll}{4}
\newcommand{\curyRec}{0.2}
\newcommand{\curyPoll}{2}
\newcommand{\renanRec}{8}
\newcommand{\renanPoll}{4}
\newcommand{\novoRec}{13}
\newcommand{\novoSeats}{1.4}
\newcommand{\zemaRecEq}{26}
\newcommand{\novoRecEq}{17}
\newcommand{\zemaRecAdv}{17}
\newcommand{\novoRecAdv}{13}
\newcommand{\lulaRecAdv}{18}
\newcommand{\flavioRecAdv}{27}
\newcommand{\caiadoRecAdv}{15}
\newcommand{\plRecAdv}{22}
\newcommand{\ptRecAdv}{9}
\newcommand{\lulaRecEq}{20}
\newcommand{\flavioRecEq}{21}
\newcommand{\plRec}{17}
\newcommand{\plSeats}{24}
\newcommand{\ptRec}{22}
\newcommand{\ptSeats}{15}
\newcommand{\leftPT}{46}
\newcommand{\leftPSOL}{16}
\newcommand{\leftPCdoB}{6}
\newcommand{\rightPL}{33}
\newcommand{\rightNOVO}{21}
\newcommand{\centreNOVO}{36}
\newcommand{\webStablePromptsPeople}{36}
\newcommand{\webMedTopPerson}{97}
\newcommand{\webMedDistinctPeople}{6}
\newcommand{\webStablePromptsParties}{31}
\newcommand{\webMedTopParty}{99}
\newcommand{\webMedDistinctParties}{5}
\newcommand{\webCiteShare}{60}
\newcommand{\webSearchShare}{32}

\newcommand{\domainsWithout}{tse.jus.br (40), www1.folha.uol.com.br (13), noticias.uol.com.br (13), divulgacandcontas.tse.jus.br (9), dadosabertos.tse.jus.br (5), gazetadopovo.com.br (3), tre-sp.jus.br (2), jota.info (2)}
\newcommand{\domainsWith}{www1.folha.uol.com.br (26), noticias.uol.com.br (12), tse.jus.br (9), gazetadopovo.com.br (7), economia.uol.com.br (4), cnnbrasil.com.br (3), camara.leg.br (3), jota.info (2)}

\setlist{nosep}

\title{Do AI chatbots tell Brazilian voters whom to vote for?\\[4pt]
\large Results outline and menu of exhibits, 2026 general election}
\author{Luca Moreno-Louzada \and Leticia Auriemo \and Andrew B.\ Hall}
\date{Working outline, 9 September 2026}

\begin{document}
\maketitle

\begin{abstract}
\noindent We ask ten API-accessible language models and the consumer ChatGPT interface whom a São Paulo voter should support for president and for federal deputy in October 2026, varying how much the voter reveals about themselves. Every system declines to advise a voter who reveals nothing. Once the voter states one policy priority, four of the ten API models name candidates or parties in most answers; adding age, gender, income and education on top changes little. Advice concentrates on parties: profiles on the left are sent to the PT, business and centre profiles to Novo, conservatives to the PL, with unanimity across models for several profiles; at the candidate level the agreement holds for president and dissolves for federal deputies. Recommendations follow the voter's stated side but not the field: Novo receives a quarter of presidential recommendations against 2\% in the polls taken during the collection window, and 17\% of São Paulo deputy recommendations against one seat in 70. The ChatGPT web interface, the surface most Brazilians use, advises in 7.5\% of prompts against 52\% for GPT-4o through the API; when it does advise on a prompt, it names the same option in 97\% of captures, and its advice-giving answers cite the press where its refusals cite the electoral court. Labels are automated (GPT-5 mini over a written codebook) and not yet human-validated; 1.6\% of responses could not be labelled, and the information contrasts are shown to hold even when those are all counted as advice or all as none.
\end{abstract}

\noindent\textbf{How to read this document.} It is a menu, not a draft. Sections 2--8 are candidate results with one exhibit each; Section 9 lists further exhibits already produced; Appendix C reproduces ten real exchanges, verbatim, with their record identifiers. Every number is read from the pipeline's output tables, never retyped.

\section{Question and design}

\paragraph{Question.} When a voter asks an AI assistant whom to vote for, does the assistant answer, and if so, whom does it name? Two follow-ups: which information about the voter unlocks advice, the person or the policy view; and does the advice converge on the same options across systems.

\paragraph{Prompts.} 64 voter bodies $\times$ 4 questions $=256$ prompts, in Portuguese. Bodies vary what the voter reveals, in five conditions: nothing (the bare question); a biography only (age, gender, income, education, state); one policy priority only (nine issue statements, one per archetype); biography plus the priority; and, on top of those, attitudes and identities, including past vote and party sympathy. We call the last the full profile. The nine archetypes come from Neto's (2024) segmentation of the Brazilian electorate: three lean left (left activist, progressive, state-dependent), two centre (social liberal, solo entrepreneur), four right (Christian conservative, agribusiness, business owner, far right). Questions cross office (president; federal deputy for São Paulo) with wording (``Em quem devo votar'' vs.\ ``Em qual candidato específico devo votar''). Appendix C prints prompts in full.

\paragraph{Systems.} Ten API configurations (GPT-4o, GPT-5.6 Luna and Sol, Claude Sonnet 5 and Opus 5, Gemini 3.1 Pro, DeepSeek V4 Pro, Grok 4.6, Llama 4 Maverick, Sabiá 4), five repetitions per prompt, collected 15--22 August 2026, after candidate registration. Plus the ChatGPT web interface, collected by Ranqia on 21--29 August with roughly 90,000 captures; 100 per prompt and day enter the coded sample. The web surface runs an unidentified model and is reported as an eleventh row, not as one of the ten.

\paragraph{Measurement.} Each answer is coded against a written codebook into independent flags: refusal language, explicit endorsement, personalized matching (directing this voter to named options because of stated priorities), negative steering, procedural guidance, and the named entities with their stance. ``Positive advice'' is endorsement or matching. The coder is GPT-5 mini run by Ranqia over the frozen 103,934-response sample; 1.6\% of responses failed the coder's evidence checks and carry no label (Table~\ref{tab:models}, last column; up to 13\% for Sabiá). A cue-based regex coder runs alongside as a second instrument (Appendix A). No human validation has been completed; the 628-response validation sheet is ready but blank.

\paragraph{Estimands and weights.} The main analysis uses the ``specific candidate'' wording, the closer analogue of the forced choice in the Japan study; the open wording enters through one wording exhibit (Result~3), and the party-level exhibits for full profiles (Results 4--6) pool both wordings because advice-giving answers are scarce for the cautious configurations. Collection days are averaged within a prompt, prompts within a configuration, configurations within any pooled number. Rates are shares of labelled responses, with 95\% intervals from the standard error across prompt cells within a configuration; unlabelled responses are reported as ``not coded''. Contrasts carry 95\% intervals from the across-pair standard error, and Appendix D reports the same estimands as per-configuration linear probability models with errors clustered by prompt. The information contrasts in Table~\ref{tab:ladder} additionally carry worst-case bounds that count the unlabelled responses all as advice or all as none; elsewhere the bounds are omitted because for eight of eleven configurations they are narrower than the intervals. These are averages over designed prompts, not prevalence among Brazilian voters.

\section{Result 1. Declining and advising are not opposites}

\begin{figure}[htbp]\centering
\includegraphics[width=.8\textwidth]{01b_advice_with_missing_bounds.pdf}
\caption{Probability of positive advice by configuration. Point: labelled responses; whiskers: 95\% interval across prompt cells. Label: share not coded.}
\label{fig:advice}
\end{figure}

\begin{table}[htbp]\centering\footnotesize
\caption{Behaviour by configuration, specific-candidate wording, 128 prompts (\%, equal prompts). Flags are independent and can overlap.}
\label{tab:models}
\begin{tabular}{lrrrrrrrr}
\toprule
 & \multicolumn{3}{c}{Positive advice} & Refusal & Disclaimer & Explicit & Procedural & Not \\
Configuration & \% & (s.e.) & 95\% interval & language & then advice & endorsement & guidance & coded \\
\midrule
GPT-4o & 51.8 & (3.7) & 44.5--59.1 & 34.1 & 0.8 & 23.2 & 78.9 & 1.3 \\
GPT-5.6 Luna & 1.5 & (0.8) & 0.0--3.0 & 100.0 & 1.5 & 0.6 & 100.0 & 0.7 \\
GPT-5.6 Sol & 0.5 & (0.3) & 0.0--1.0 & 100.0 & 0.5 & 0.1 & 100.0 & 0.0 \\
Claude Sonnet 5 & 5.4 & (1.1) & 3.2--7.6 & 100.0 & 5.4 & 0.1 & 99.7 & 2.7 \\
Claude Opus 5 & 4.3 & (1.1) & 2.2--6.5 & 100.0 & 4.3 & 0.0 & 100.0 & 2.2 \\
Gemini 3.1 Pro & 28.1 & (3.2) & 21.9--34.3 & 100.0 & 28.1 & 1.4 & 100.0 & 2.9 \\
DeepSeek V4 Pro & 52.1 & (3.8) & 44.7--59.5 & 75.5 & 27.6 & 32.4 & 95.5 & 7.3 \\
Grok 4.6 & 60.1 & (3.9) & 52.5--67.7 & 67.3 & 27.4 & 42.0 & 96.5 & 5.9 \\
Llama 4 Maverick & 10.9 & (1.5) & 8.0--13.7 & 24.1 & 2.3 & 4.0 & 91.3 & 1.9 \\
Sabi\'a 4 & 43.3 & (3.3) & 36.7--49.8 & 74.4 & 23.1 & 22.7 & 99.0 & 12.2 \\
ChatGPT (web) & 7.7 & (1.5) & 4.8--10.6 & 99.9 & 7.7 & 0.8 & 99.5 & 0.7 \\
\bottomrule
\end{tabular}

\end{table}

Three groups (Figure~\ref{fig:advice}, Table~\ref{tab:models}). Grok, GPT-4o, DeepSeek and Sabiá advise in 43--60\% of prompts and endorse explicitly in 23--42\%. Gemini advises in 28\% but almost never endorses: it matches options to the voter after a disclaimer. The OpenAI 5.6 models, both Claude models, Llama and the ChatGPT web surface advise in 0.5--11\%. Refusal language is universal for the cautious group and for Gemini, and still present in two thirds to three quarters of Grok, DeepSeek and Sabiá answers; only GPT-4o usually skips the disclaimer. ``I cannot tell you whom to vote for'' followed by a shortlist is the modal advice-giving answer for every configuration except GPT-4o (Appendix C, examples 4 and 9).

\section{Result 2. The issue unlocks advice; the person does not}

\begin{figure}[htbp]\centering
\includegraphics[width=\textwidth]{03_information_ladder_levels.pdf}
\caption{Advice probability by what the voter reveals, by configuration. Biography and issue are separate additions to the bare question; the last two conditions add the issue to the biography, then attitudes and identities. Band: 95\% interval across prompt cells.}
\label{fig:ladder}
\end{figure}

\begin{table}[htbp]\centering\footnotesize\setlength{\tabcolsep}{3pt}
\caption{Advice probability by what the voter reveals, with the two within-profile contrasts. Bounds count unlabelled responses all as advice or all as no advice. $\dagger$: sign not robust to the bounds.}
\label{tab:ladder}
\begin{tabular}{lrrrrrrrrrr}
\toprule
 & \multicolumn{5}{c}{Advice probability by what the voter reveals (\%)} & \multicolumn{3}{c}{Issue added to biography} & \multicolumn{2}{c}{Attitudes added} \\
\cmidrule(lr){2-6} \cmidrule(lr){7-9} \cmidrule(lr){10-11}
Configuration & Nothing & Biography & Issue & Bio + issue & + Attitudes & pp & (s.e.) & bounds & pp & (s.e.) \\
\midrule
GPT-4o & 0 & 2 & 74 & 64 & 81 & 62 & (5.6) & 62--62 & 18 & (5.2) \\
GPT-5.6 Luna & 0 & 0 & 4 & 1 & 3 & 1 & (0.6) & 1--1 & 2 & (1.5) \\
GPT-5.6 Sol & 0 & 0 & 3 & 0 & 0 & 0 & (0.3) & 0--0 & 0 & (0.3) \\
Claude Sonnet 5 & 0 & 0 & 8 & 3 & 12 & 3 & (1.6) & 2--3 & 8 & (3.4) \\
Claude Opus 5 & 0 & 0 & 7 & 5 & 6 & 5 & (2.3) & 3--8 & 1 & (2.5)$^{\dagger}$ \\
Gemini 3.1 Pro & 0 & 0 & 19 & 28 & 62 & 28 & (5.5) & 27--29 & 34 & (6.4) \\
DeepSeek V4 Pro & 0 & 0 & 77 & 56 & 90 & 56 & (5.3) & 52--58 & 34 & (5.5) \\
Grok 4.6 & 0 & 1 & 81 & 77 & 94 & 76 & (4.3) & 69--78 & 17 & (4.5) \\
Llama 4 Maverick & 0 & 1 & 8 & 11 & 23 & 9 & (2.0) & 9--11 & 13 & (3.5) \\
Sabi\'a 4 & 0 & 7 & 39 & 48 & 80 & 41 & (5.2) & 32--46 & 31 & (6.6) \\
ChatGPT (web) & 0 & 0 & 2 & 14 & 13 & 14 & (3.6) & 13--15 & -1 & (1.8)$^{\dagger}$ \\
\bottomrule
\end{tabular}

\end{table}

Pooling configurations equally, advice is given in \poolLone\% of bare prompts, \poolLthree\% of biography-only prompts, \poolLtwo\% of issue-only prompts, \poolLfour\% with biography and issue, and \poolLfive\% with attitudes added. The pattern is the same in every configuration (Figure~\ref{fig:ladder}): the bare question and the biography sit at zero, the issue statement alone produces most of the advice, and adding the biography to the issue does not raise it (pooled $-2$ pp). Adding the issue to a biography raises advice by 62 pp for GPT-4o and 69--78 pp for Grok at the worst-case bounds; the sign is robust for all eleven configurations (Table~\ref{tab:ladder}). Adding attitudes and identities on top adds a further 15 pp pooled, but the sign is not robust to the bounds for Claude Opus and the web surface. This is the Miyazaki--Hall (2026) result for Japan: what the voter says about policy, not who they are, is what the system acts on. Appendix C, examples 2 and 3, show the same model refusing a biography and endorsing a candidate for an issue.

\begin{figure}[htbp]\centering
\includegraphics[width=\textwidth]{03c_ladder_contrasts_ci.pdf}
\caption{The four information contrasts as matched mean differences with 95\% intervals (across-pair standard errors), by configuration.}
\label{fig:ladderci}
\end{figure}

\begin{figure}[htbp]\centering
\includegraphics[width=.9\textwidth]{13_issue_heterogeneity.pdf}
\caption{Which stated priorities unlock advice: the change when the issue is added to the biography, by profile and configuration. Bar: equal-model mean.}
\label{fig:heterogeneity}
\end{figure}

The step differs by issue (Figure~\ref{fig:heterogeneity}). Business-owner and far-right statements (tax cuts and privatization; order and nationalism) unlock advice most, 38--48 pp on average; the state-dependent and social-liberal statements (cost of living and social programmes; democracy and depolarization) least, 11--12 pp. The ranking of configurations is the same within every issue.

\section{Result 3. Wording matters at the margin; office and gender do not}

\begin{figure}[htbp]\centering
\includegraphics[width=\textwidth]{05b_wording_levels.pdf}
\caption{Open question versus specific-candidate question, same prompts otherwise. Main tables use the specific wording.}
\label{fig:wording}
\end{figure}

The open wording (``Em quem devo votar'') yields 5 pp more advice pooled than the specific-candidate wording used in the main tables, and 3 pp less refusal language, with the same ordering of configurations and the same ladder shape (Figure~\ref{fig:wording}). The difference sits with the cautious configurations, which are not quite silent when asked openly: GPT-5.6 Luna 10\% vs.\ 2\%, GPT-5.6 Sol 6\% vs.\ 0.5\%, Claude Sonnet 11\% vs.\ 5\%, the web surface 13\% vs.\ 8\%; the four advising configurations move by 0--9 pp. Moving from president to federal deputy changes advice by 0 pp pooled, with heterogeneity: GPT-4o and Grok advise more down-ballot ($+12$, $+10$), DeepSeek and the web surface less ($-13$, $-11$). Down-ballot advice is more often a party than a person for Gemini and Sabiá (27\% and 24\% party-only; Figure~\ref{fig:office}). A woman rather than a man with the same profile receives advice 2 pp less often (Figure~\ref{fig:gender}).

\section{Result 4. Advice converges on parties; on candidates only where the field is small}

\begin{figure}[htbp]\centering
\includegraphics[width=\textwidth]{11_convergence_top_party.pdf}
\caption{Share of recommended parties going to the top party, by configuration and full profile, both wordings pooled. Only configuration--profile pairs with any advice appear.}
\label{fig:convergence}
\end{figure}

\begin{table}[htbp]\centering\footnotesize\setlength{\tabcolsep}{3pt}
\caption{Convergence at the party level, full profiles, conditional on advice. ``Agreeing'' is the share of advice-giving configurations whose top party is the modal one. Effective parties $=1/\mathrm{HHI}$.}
\label{tab:convergence}
\begin{tabular}{llrlrrr}
\toprule
Profile & Office & Models & Modal top party & Agreeing (\%) & Top share (\%) & Eff.\ parties \\
\midrule
Left activist & President & 9 & PT & 100 & 95 & 1.1 \\
Left activist & Deputy & 11 & PT & 100 & 74 & 1.5 \\
Progressive & President & 9 & PT & 78 & 48 & 2.6 \\
Progressive & Deputy & 9 & PSOL & 67 & 38 & 3.6 \\
State-dependent & President & 7 & PT & 57 & 50 & 2.4 \\
State-dependent & Deputy & 5 & Missao & 20 & 43 & 2.8 \\
Social liberal & President & 10 & NOVO & 60 & 45 & 2.9 \\
Social liberal & Deputy & 8 & NOVO & 62 & 24 & 5.6 \\
Solo entrepreneur & President & 10 & NOVO & 60 & 39 & 3.2 \\
Solo entrepreneur & Deputy & 9 & NOVO & 78 & 48 & 2.7 \\
Christian conservative & President & 9 & PL & 78 & 42 & 3.1 \\
Christian conservative & Deputy & 8 & PL & 100 & 36 & 3.7 \\
Agribusiness & President & 9 & PL & 89 & 42 & 3.0 \\
Agribusiness & Deputy & 8 & PL & 100 & 41 & 3.4 \\
Business owner & President & 11 & NOVO & 100 & 60 & 2.1 \\
Business owner & Deputy & 9 & NOVO & 100 & 46 & 2.9 \\
Far right & President & 10 & PL & 80 & 52 & 2.4 \\
Far right & Deputy & 9 & PL & 78 & 44 & 3.0 \\
\bottomrule
\end{tabular}

\end{table}

\begin{table}[htbp]\centering\footnotesize\setlength{\tabcolsep}{3pt}
\caption{The same measures at the candidate level. Effective candidates $=1/\mathrm{HHI}$ over named people.}
\label{tab:convergenceperson}
\begin{tabular}{llrlrrr}
\toprule
Profile & Office & Models & Modal top candidate & Agreeing (\%) & Top share (\%) & Eff.\ candidates \\
\midrule
Left activist & President & 9 & Luiz Inacio Lula Da Silva & 100 & 91 & 1.2 \\
Progressive & President & 8 & Luiz Inacio Lula Da Silva & 75 & 57 & 2.1 \\
State-dependent & President & 7 & Luiz Inacio Lula Da Silva & 57 & 53 & 2.2 \\
Social liberal & President & 10 & Romeu Zema & 80 & 46 & 3.0 \\
Solo entrepreneur & President & 10 & Romeu Zema & 60 & 40 & 3.2 \\
Christian conservative & President & 9 & Flavio Bolsonaro & 78 & 41 & 3.2 \\
Agribusiness & President & 9 & Flavio Bolsonaro & 78 & 41 & 3.1 \\
Business owner & President & 11 & Romeu Zema & 91 & 60 & 2.1 \\
Far right & President & 10 & Flavio Bolsonaro & 90 & 48 & 2.6 \\
Left activist & Deputy & 8 & Juliana Cardoso & 25 & 17 & 9.1 \\
Progressive & Deputy & 6 & Erika Hilton & 50 & 33 & 4.7 \\
State-dependent & Deputy & 5 & Alencar Santana & 20 & 44 & 2.6 \\
Social liberal & Deputy & 10 & Kim Kataguiri & 40 & 45 & 2.5 \\
Solo entrepreneur & Deputy & 7 & Adriana Ventura & 29 & 31 & 4.6 \\
Christian conservative & Deputy & 7 & Guilherme Derrite & 29 & 22 & 6.8 \\
Agribusiness & Deputy & 7 & Ricardo Salles & 29 & 19 & 9.8 \\
Business owner & Deputy & 6 & Marina Helena & 50 & 26 & 5.9 \\
Far right & Deputy & 7 & Renato Bolsonaro & 43 & 34 & 4.0 \\
\bottomrule
\end{tabular}

\end{table}

Conditional on giving advice, and pooling both wordings, configurations send the same profile to the same party (Table~\ref{tab:convergence}). Every advice-giving configuration sends the left activist to the PT (92\% of presidential recommendations, effective number of parties 1.1), the business owner to Novo, and the Christian conservative to the PL for deputy. Agreement is weaker for the state-dependent profile, whose stated priorities (cost of living, social programmes) are read as PT by some systems and as centre-right by others, and for the social liberal, where deputy recommendations spread over six effective parties.

The claim has to be stated at the right level. At the candidate level (Table~\ref{tab:convergenceperson}) convergence holds for president, where the field is small: the modal candidate is shared by 57--100\% of advice-giving configurations, with the left activist sent to Lula (91\% of named people), the business owner to Zema (62\%) and the far right to Flávio Bolsonaro (48\%). But the right and centre presidential profiles are split among Flávio, Zema and Caiado, about three effective candidates, which is what the person heatmap in Figure~\ref{fig:persons} shows: the same profile lands on different names in different systems, and shortlists spread mass across several. For federal deputy the party consensus does not carry to names at all: 4--9 effective candidates per profile and the modal name shared by only 20--67\% of configurations. Convergence is about the profile's stated side and its flagship party, not about a specific person, and not about the model.

\section{Result 5. Same position on the scale, different treatment}

\begin{figure}[htbp]\centering
\includegraphics[width=\textwidth]{10_party_share_by_position.pdf}
\caption{Share of recommended options by party, against the party's Zucco--Power left--right position, by profile side and office. Full profiles, conditional on advice; profiles then configurations weighted equally.}
\label{fig:position}
\end{figure}

If systems matched voters to positions, parties at the same position would be recommended at similar rates. They are not (Figure~\ref{fig:position}). Among left profiles asking about federal deputies, the PT takes \leftPT\% of recommended parties; PSOL, further left on the scale, takes \leftPSOL\%, and PCdoB, at the same position as PSOL, \leftPCdoB\%. PDT and PSB, closer to the centre, sit at 5\% each. Among right profiles, the PL takes \rightPL\% and Novo \rightNOVO\%; PP, Republicanos and União, at the same positions as the PL and Novo, take 8--14\% each. Among centre profiles, Novo takes \centreNOVO\% of deputy recommendations, more than every party actually near the centre of the scale combined. This is the Japanese Communist Party pattern: the system picks a salient flagship for each side, not the nearest party.

\section{Result 6. Recommendations do not track the field}

\begin{figure}[htbp]\centering
\includegraphics[width=\textwidth]{12_recommendations_vs_benchmarks.pdf}
\caption{Share of recommendations next to the electoral benchmark: Genial/Quaest first-round poll of 10--13 August 2026 for president; 2022 seat shares for São Paulo federal deputies. Full profiles, conditional on advice, both wordings. Three weightings: models equal and profiles by population share (navy); share of all advice given, so models count in proportion to how often they advise (teal); models and profiles equal (grey).}
\label{fig:benchmarks}
\end{figure}

The nine profiles are not an electorate, so for this comparison they are weighted by their share of the population in Neto's segmentation (Christian conservatives 27\%, state-dependent 23\%, agribusiness 13\%, progressives 11\%, left activists 7\%, business owners 6\%, social liberals 5\%, solo entrepreneurs 5\%, far right 3\%). Two model weightings are shown, because they answer different questions. With models weighted equally (what systems say when they advise), Romeu Zema (Novo) receives \zemaRec\% of presidential recommendations (\zemaRecEq\% with equal profile weights) and polls at \zemaPoll\%; Lula receives \lulaRec\% against \lulaPoll\%, Flávio Bolsonaro \flavioRec\% against \flavioPoll\%, Caiado \caiadoRec\% against \caiadoPoll\%. Weighting models by how often they advise (the composition of all advice actually given, dominated by Grok, GPT-4o, DeepSeek and Sabiá) moves the picture toward the right's frontrunner: Flávio Bolsonaro \flavioRecAdv\%, Lula \lulaRecAdv\%, Zema \zemaRecAdv\%, Caiado \caiadoRecAdv\% (Figure~\ref{fig:benchmarks}). Under either weighting Zema is recommended about ten times as often as he polls, the two third-way candidates take roughly a third of recommendations against 6\% in the poll, and Lula trails his poll share by 15--21 points. For São Paulo deputies, Novo receives \novoRec\% of party recommendations under both weightings (\novoRecEq\% equal-weighted) with \novoSeats\% of seats; the PL \plRec\% or \plRecAdv\% against \plSeats\%, the PT \ptRec\% or \ptRecAdv\% against \ptSeats\%. The bare question cannot carry this comparison, because almost no configuration advises on it. The comparison remains descriptive: the segments' shares come from a 2024 cluster model, each segment is one synthetic voter, and the poll measures intentions, not what a voter would be told.

\section{Result 7. The consumer surface}

Most Brazilians who ask a chatbot about the election ask ChatGPT in a browser or app, not an API. The web arm is therefore the exposure that matters, and it behaves differently from every API configuration of the same company.

\subsection{Same company, different surface}

\begin{figure}[htbp]\centering
\includegraphics[width=.75\textwidth]{14_web_vs_api_ladder.pdf}
\caption{Advice probability by what the voter reveals: OpenAI API configurations and the ChatGPT web surface.}
\label{fig:webladder}
\end{figure}

\begin{table}[htbp]\centering\small
\caption{Advice probability (\%) by what the voter reveals, OpenAI configurations and the web surface, specific-candidate wording.}
\label{tab:webladder}
\begin{tabular}{lrrrrr}
\toprule
Configuration & Nothing & Biography & Issue & Bio + issue & + Attitudes \\
\midrule
GPT-4o & 0 & 2 & 74 & 64 & 81 \\
GPT-5.6 Luna & 0 & 0 & 4 & 1 & 3 \\
GPT-5.6 Sol & 0 & 0 & 3 & 0 & 0 \\
ChatGPT (web) & 0 & 0 & 2 & 14 & 13 \\
\bottomrule
\end{tabular}

\end{table}

ChatGPT on the web advises in \webAdvice\% of prompts against 52\% for GPT-4o through the API and 0.5--1.5\% for the GPT-5.6 configurations (Figure~\ref{fig:webladder}, Table~\ref{tab:webladder}). It never advises on a bare or biography-only prompt, and, unlike the API models, barely reacts to the issue alone (2\%): it needs the biography and the issue together (14\%) or the full profile (13\%). Refusal language is universal on the web (100\%), against 34\% for GPT-4o. Every advice-giving web answer opens with a refusal (Appendix C, example 9).

\subsection{How predictable is the web answer?}

\begin{figure}[htbp]\centering
\includegraphics[width=.75\textwidth]{08_web_advice_probabilities.pdf}
\caption{Estimated advice probability for each of the 128 specific-candidate prompts on the web surface, averaging collection days equally.}
\label{fig:web}
\end{figure}

\begin{figure}[htbp]\centering
\includegraphics[width=.9\textwidth]{15_web_option_stability.pdf}
\caption{What the web surface does with each full-profile prompt, 300--400 labelled captures per bar: no advice (grey), advice naming the prompt's lead option (navy), advice naming something else (coral). The lead option is the person, or the party where no person is named, that advice-giving captures name most often; the text gives it and the runner-up with their shares of the advice-giving captures. Both shares are high when the answer is a recurring shortlist.}
\label{fig:webstability}
\end{figure}

Across the 128 prompts, \webNeverAdvise\% never produce advice in any capture and \webMostlyAdvise\% produce it in most (Figure~\ref{fig:web}). Advice is rare but, where it occurs, it is the same advice every time (Figure~\ref{fig:webstability}). For president, every full-profile prompt has one lead name that appears in 90--100\% of its advice-giving captures: Lula for the left activist, progressive and state-dependent man; Romeu Zema for the business owner, solo entrepreneur and social liberal; Flávio Bolsonaro for the far-right, Christian-conservative and agribusiness profiles and the state-dependent woman. Where a second name appears it is a recurring shortlist, not variation: the conservative and agribusiness profiles get Flávio Bolsonaro and Ronaldo Caiado together in about 80\% of their advice-giving captures, the social liberal Zema and Caiado, the entrepreneur Zema and Flávio Bolsonaro; the left activist gets Lula alone. How often the web advises varies far more than what it advises: from 77--79\% of captures for the business owner and 45\% for the solo entrepreneur to 2--4\% for the progressive and state-dependent profiles. For deputies advice is rare (0--33 captures per prompt), and where it exists it is the PT as a party for the left activist, Marina Helena (Novo) for the solo entrepreneur, and Novo or a Novo candidate for the business owner. Prompts with the biography and the issue but no attitudes give the same picture, with one telling exception: the agribusiness man is sent to Ronaldo Caiado (96\%) before attitudes are added and to Flávio Bolsonaro (95\%) after, once the profile states past support for Jair Bolsonaro.

\subsection{Where the web surface sends each profile}

\begin{table}[htbp]\centering\footnotesize
\caption{Full profiles, both wordings pooled. Web advice rate, its top recommended party and that party's share of web recommendations, next to the modal top party across API configurations and the share of them agreeing.}
\label{tab:webparties}
\begin{tabular}{llrlrlr}
\toprule
 & & \multicolumn{3}{c}{ChatGPT web} & \multicolumn{2}{c}{API configurations} \\
\cmidrule(lr){3-5} \cmidrule(lr){6-7}
Profile & Office & advice (\%) & top party & share (\%) & modal top party & agreeing (\%) \\
\midrule
Left activist & President & 43 & PT & 100 & PT & 100 \\
Progressive & President & 5 & PT & 51 & PT & 75 \\
State-dependent & President & 8 & PT & 48 & PT & 50 \\
Social liberal & President & 20 & PSD & 33 & NOVO & 67 \\
Solo entrepreneur & President & 52 & NOVO & 41 & NOVO & 56 \\
Christian conservative & President & 21 & PL & 32 & PL & 75 \\
Agribusiness & President & 15 & PL & 44 & PL & 88 \\
Business owner & President & 81 & NOVO & 54 & NOVO & 100 \\
Far right & President & 18 & PL & 37 & PL & 78 \\
Left activist & Deputy & 7 & PT & 85 & PT & 100 \\
Progressive & Deputy & 1 & PSOL & 47 & PSOL & 62 \\
State-dependent & Deputy & 0 & NOVO & 100 & MISSAO & 25 \\
Social liberal & Deputy & 3 & NOVO & 24 & NOVO & 57 \\
Solo entrepreneur & Deputy & 13 & NOVO & 70 & NOVO & 75 \\
Christian conservative & Deputy & 1 & PL & 44 & PL & 100 \\
Agribusiness & Deputy & 1 & PL & 35 & PL & 100 \\
Business owner & Deputy & 14 & NOVO & 71 & NOVO & 100 \\
Far right & Deputy & 1 & PL & 44 & PL & 75 \\
\bottomrule
\end{tabular}

\end{table}

The web surface's advice is concentrated in two profiles (Table~\ref{tab:webparties}): the business owner (81\% of presidential prompts receive advice, 54\% of it Novo) and the solo entrepreneur (52\%, Novo), followed by the left activist (43\%, all PT). Conservative, agribusiness and far-right profiles receive advice in 15--21\% of presidential prompts, progressive and state-dependent in 5--8\%. Where it advises, the web surface picks the same party as the API modal answer in every profile except the social liberal (PSD rather than Novo).

\subsection{What the web surface cites}

\begin{figure}[htbp]\centering
\includegraphics[width=.9\textwidth]{16_web_citation_sources.pdf}
\caption{Sources cited in visible web answers, by whether the answer gives advice; one link occurrence is one citation.}
\label{fig:webcites}
\end{figure}

\begin{table}[htbp]\centering\small
\caption{Share of cited links by source type (\%), web answers with and without advice.}
\label{tab:websources}
\begin{tabular}{lrr}
\toprule
Source type & Answers without advice & Answers with advice \\
\midrule
Press and news & 35.3 & 66.6 \\
Official (TSE, Congress, government) & 58.6 & 17.2 \\
Other & 5.0 & 11.6 \\
Party and movement sites & 1.0 & 3.6 \\
Social media & 0.1 & 0.6 \\
Wikipedia & 0.1 & 0.4 \\
\bottomrule
\end{tabular}

\end{table}

\webCiteShare\% of web answers display at least one citation and \webSearchShare\% ran a visible search. The sources differ sharply by what the answer does (Figure~\ref{fig:webcites}, Table~\ref{tab:websources}). Answers without advice cite official sources first: \domainsWithout. Answers with advice cite the press: \domainsWith. The electoral court's own site falls from 40\% of citations to 9\% when the answer names a candidate. This is the mechanism the Japan paper documents: advice comes from the retrievable news environment, refusal from the official one.

\section{Further exhibits produced}

\begin{figure}[htbp]\centering
\includegraphics[width=.9\textwidth]{02_response_categories.pdf}
\caption{Mutually exclusive response categories by configuration, specific-candidate wording; ``not coded'' is the unlabelled share.}
\label{fig:categories}
\end{figure}

\begin{figure}[htbp]\centering
\includegraphics[width=\textwidth]{04_office_and_specificity.pdf}
\caption{Form of positive advice by office: one named person, several, or a party without a person.}
\label{fig:office}
\end{figure}

\begin{figure}[htbp]\centering
\includegraphics[width=.9\textwidth]{05_wording_and_gender.pdf}
\caption{Matched mean differences in advice probability with 95\% intervals: specific-candidate minus open wording, and woman minus man.}
\label{fig:gender}
\end{figure}

\begin{figure}[htbp]\centering
\includegraphics[width=\textwidth]{06_profiles_receiving_advice.pdf}
\caption{Advice probability by full profile and office, five configurations.}
\label{fig:profiles}
\end{figure}

\begin{figure}[htbp]\centering
\includegraphics[width=\textwidth]{07_person_recommendations_president.pdf}
\caption{Presidential candidates recommended, by full profile, five configurations; share of all full-profile responses.}
\label{fig:persons}
\end{figure}

\begin{figure}[htbp]\centering
\includegraphics[width=.75\textwidth]{09_web_candidate_consistency.pdf}
\caption{Web surface: probability that two same-day advice-giving captures of a prompt name the identical set of people.}
\label{fig:webpairs}
\end{figure}

\begin{table}[htbp]\centering\footnotesize
\caption{Secondary flags by configuration (\%, specific-candidate wording).}
\label{tab:flags}
\begin{tabular}{lrrrrrrr}
\toprule
 & \multicolumn{3}{c}{Form of positive advice} & Negative & Names any & Recites a & Claims 2026 \\
Configuration & one person & several & party only & steering & person & slate & undefined \\
\midrule
GPT-4o & 14.1 & 36.7 & 0.9 & 0.8 & 63.7 & 9.1 & 15.2 \\
GPT-5.6 Luna & 0.7 & 0.3 & 0.5 & 0.3 & 2.9 & 0.2 & 22.4 \\
GPT-5.6 Sol & 0.4 & 0.1 & 0.0 & 0.1 & 1.4 & 0.1 & 66.0 \\
Claude Sonnet 5 & 0.7 & 2.4 & 2.2 & 0.7 & 14.1 & 7.1 & 29.3 \\
Claude Opus 5 & 1.7 & 1.4 & 1.2 & 0.7 & 52.2 & 46.7 & 21.6 \\
Gemini 3.1 Pro & 1.1 & 9.0 & 17.9 & 1.3 & 18.2 & 3.3 & 37.6 \\
DeepSeek V4 Pro & 22.2 & 28.3 & 1.6 & 12.0 & 66.3 & 14.2 & 56.5 \\
Grok 4.6 & 39.9 & 17.8 & 2.4 & 18.0 & 79.1 & 19.1 & 18.6 \\
Llama 4 Maverick & 2.6 & 7.7 & 0.5 & 0.2 & 22.3 & 7.6 & 55.7 \\
Sabi\'a 4 & 9.9 & 19.9 & 13.4 & 13.8 & 56.1 & 14.9 & 63.4 \\
ChatGPT (web) & 3.5 & 4.0 & 0.2 & 1.5 & 38.9 & 19.2 & 14.7 \\
\bottomrule
\end{tabular}

\end{table}

\begin{table}[htbp]\centering\footnotesize\setlength{\tabcolsep}{3pt}
\caption{Ground given for the refusal, first explicit ground in the answer (\%, specific-candidate wording).}
\label{tab:grounds}
\begin{tabular}{lrrrrr}
\toprule
Configuration & No refusal & Neutrality, autonomy & Candidacies undefined & Knowledge limits & Unclear ground \\
\midrule
GPT-4o & 65.9 & 25.9 & 5.2 & 1.9 & 1.1 \\
GPT-5.6 Luna & 0.0 & 99.4 & 0.3 & 0.0 & 0.3 \\
GPT-5.6 Sol & 0.0 & 98.8 & 0.9 & 0.0 & 0.3 \\
Claude Sonnet 5 & 0.0 & 96.5 & 1.7 & 1.8 & 0.0 \\
Claude Opus 5 & 0.0 & 99.8 & 0.0 & 0.2 & 0.0 \\
Gemini 3.1 Pro & 0.0 & 100.0 & 0.0 & 0.0 & 0.0 \\
DeepSeek V4 Pro & 24.5 & 70.0 & 5.5 & 0.0 & 0.0 \\
Grok 4.6 & 32.7 & 67.2 & 0.1 & 0.0 & 0.0 \\
Llama 4 Maverick & 75.9 & 3.1 & 14.9 & 6.0 & 0.2 \\
Sabi\'a 4 & 25.6 & 66.2 & 7.7 & 0.3 & 0.1 \\
ChatGPT (web) & 0.1 & 99.5 & 0.1 & 0.0 & 0.3 \\
\bottomrule
\end{tabular}

\end{table}

\begin{table}[htbp]\centering\small
\caption{Mean Zucco--Power position of recommended parties ($-1$ left to $+1$ right), by full profile; equal weight across configurations with a defined value. Both wordings pooled.}
\label{tab:ideology}
\begin{tabular}{lrrrr}
\toprule
 & \multicolumn{2}{c}{President} & \multicolumn{2}{c}{Federal deputy} \\
\cmidrule(lr){2-3} \cmidrule(lr){4-5}
Profile & mean $\zeta$ & models & mean $\zeta$ & models \\
\midrule
Left activist & -0.69 & 1 & -0.68 & 1 \\
Progressive & -0.67 & 1 & -0.68 & 1 \\
State-dependent & -0.08 & 1 & 0.10 & 1 \\
Social liberal & 0.50 & 1 & 0.38 & 1 \\
Solo entrepreneur & 0.58 & 1 & 0.62 & 1 \\
Christian conservative & 0.47 & 1 & 0.51 & 1 \\
Agribusiness & 0.47 & 1 & 0.53 & 1 \\
Business owner & 0.66 & 1 & 0.62 & 1 \\
Far right & 0.51 & 1 & 0.54 & 1 \\
\bottomrule
\end{tabular}

\end{table}

Reading notes. Figure~\ref{fig:categories}: for the cautious configurations, the modal answer is procedural guidance, not a bare refusal; Claude Opus is the one system whose modal answer is neutral information about the field. Table~\ref{tab:flags}: only Grok, DeepSeek and Sabiá steer voters away from options (12--18\%); slate recitation (four or more names, no advice) is the Claude Opus signature. Table~\ref{tab:grounds}: refusals are almost always grounded in neutrality or voter autonomy; ``candidacies are not yet defined'' is rare after registration, and knowledge limits are cited mainly by GPT-4o and Llama. Table~\ref{tab:ideology}: the recommended party's position tracks the profile's side, $-0.7$ for the two left profiles and $+0.5$ to $+0.7$ for the right, with the state-dependent profile the one that lands near the centre.

\section{Robustness and caveats}

\begin{enumerate}
\item \textbf{Unlabelled responses.} 1,687 of 103,934 coded responses (1.6\%) carry no label; 6--15\% for Grok, DeepSeek and Sabiá, concentrated in long, name-dense answers to the richer profiles that are more likely to contain advice (Appendix B). Results 1, 2, 4--7 survive counting them all as advice or all as no advice; the attitudes step does not for two configurations. The failures are evidence-quotation checks in the coder's output, not missing answers.
\item \textbf{Coder validity.} Labels are automated and not human-validated. The regex instrument agrees with the semantic coder on 65--97\% of responses by configuration and reproduces the configuration ordering (Appendix A); the disagreement review found the semantic coder right in most cases and the regex missing soft matching.
\item \textbf{Wording.} Pooling both wordings instead of the specific-candidate wording raises levels 2--8 pp for the cautious configurations and leaves the ladder, the ordering and the convergence pattern unchanged (Result~3).
\item \textbf{Not a voter population.} Nine designed profiles, equal weights; the bare-question baseline has 20 API draws per configuration. Zucco--Power positions are legislator-survey placements; parties without a position (Missão, Avante, PRTB) are shown but not scaled.
\item \textbf{Candidate validity is unmeasured.} No external registry has been joined; names are reported as stated, and party affiliations are those the answer itself attaches.
\item \textbf{Web arm.} Captures come from one account setup over nine days; execution dependence within a day is unknown, so no interval is placed on web rates. Collection days are averaged, never compared.
\end{enumerate}

\section{Open items}

\begin{itemize}
\item Human spot check of 15 responses (unlabelled rows the regex calls no-advice; soft-matching cases) and, later, the 628-row validation sheet.
\item Decide the blog storyline; then a reduced pipeline that produces only the chosen exhibits.
\end{itemize}

\appendix
\section{Regex and semantic coder agreement}
\begin{table}[htbp]\centering\small
\caption{Agreement between the cue-based regex coder and the semantic coder on labelled responses (\%), all prompts.}
\begin{tabular}{lrrrr}
\toprule
Configuration & Responses & Positive advice & Refusal language & Recommended-party set \\
\midrule
GPT-4o &  1,255 & 75.6 & 87.4 & 61.3 \\
GPT-5.6 Luna &  1,261 & 96.0 & 95.2 & 96.0 \\
GPT-5.6 Sol &  1,270 & 96.9 & 99.1 & 97.1 \\
Claude Sonnet 5 &  1,240 & 91.9 & 99.4 & 91.4 \\
Claude Opus 5 &  1,252 & 95.3 & 96.5 & 95.0 \\
Gemini 3.1 Pro &  1,249 & 73.6 & 97.0 & 70.1 \\
DeepSeek V4 Pro &  1,187 & 74.1 & 91.1 & 59.9 \\
Grok 4.6 &  1,176 & 65.3 & 83.8 & 56.1 \\
Llama 4 Maverick &  1,254 & 90.9 & 81.8 & 91.0 \\
Sabi\'a 4 &  1,078 & 76.6 & 78.5 & 60.9 \\
ChatGPT (web) & 90,025 & 91.9 & 97.8 & 91.8 \\
\bottomrule
\end{tabular}

\end{table}

\section{Unlabelled share by configuration and level}
\begin{table}[htbp]\centering\small
\caption{Share of responses without a semantic label (\%) by what the voter reveals, specific-candidate wording.}
\begin{tabular}{lrrrrr}
\toprule
Configuration & Nothing & Biography & Issue & Bio + issue & + Attitudes \\
\midrule
GPT-4o & 0 & 0 & 2 & 1 & 3 \\
GPT-5.6 Luna & 0 & 0 & 4 & 0 & 0 \\
GPT-5.6 Sol & 0 & 0 & 0 & 0 & 0 \\
Claude Sonnet 5 & 0 & 1 & 6 & 0 & 6 \\
Claude Opus 5 & 0 & 2 & 1 & 3 & 3 \\
Gemini 3.1 Pro & 0 & 0 & 0 & 3 & 8 \\
DeepSeek V4 Pro & 0 & 0 & 9 & 6 & 16 \\
Grok 4.6 & 0 & 0 & 6 & 9 & 9 \\
Llama 4 Maverick & 0 & 0 & 0 & 1 & 6 \\
Sabi\'a 4 & 0 & 1 & 7 & 13 & 25 \\
ChatGPT (web) & 0 & 0 & 0 & 1 & 1 \\
\bottomrule
\end{tabular}

\end{table}

\section{Regression view}
One linear probability model per configuration on labelled responses, both wordings: positive advice on level dummies (biography-only reference), office, wording and gender; standard errors clustered by prompt cell. In this balanced design the coefficients equal the matched differences of Results 2--3; the table adds the uncertainty.
\begin{table}[htbp]\centering\footnotesize\setlength{\tabcolsep}{3pt}
\caption{Coefficients in percentage points (cell-clustered standard errors), outcome positive advice.}
\begin{tabular}{lrrrrrr}
\toprule
 & \multicolumn{3}{c}{Level, vs biography only} & Deputy & Specific & Woman \\
\cmidrule(lr){2-4}
Configuration & Issue only & Bio + issue & Attitudes & vs president & vs open & vs man \\
\midrule
GPT-4o & 73 (4.9) & 66 (3.7) & 83 (2.9) & 13 (2.8) & -4 (2.9) & -1 (3.0) \\
GPT-5.6 Luna & 20 (4.2) & 3 (1.1) & 8 (2.0) & -3 (1.7) & -9 (1.7) & -2 (1.5) \\
GPT-5.6 Sol & 11 (3.2) & 5 (1.5) & 3 (1.0) & -3 (1.3) & -6 (1.3) & 0 (1.2) \\
Claude Sonnet 5 & 14 (4.0) & 4 (1.2) & 18 (2.5) & 1 (1.8) & -5 (1.9) & 0 (1.8) \\
Claude Opus 5 & 8 (3.9) & 5 (1.4) & 9 (1.8) & -2 (1.7) & -2 (1.7) & -4 (1.5) \\
Gemini 3.1 Pro & 25 (4.9) & 31 (3.8) & 69 (3.8) & 1 (3.3) & -7 (3.2) & -3 (3.5) \\
DeepSeek V4 Pro & 80 (3.9) & 55 (3.6) & 86 (3.1) & -5 (2.7) & 2 (2.7) & 1 (3.0) \\
Grok 4.6 & 80 (4.3) & 79 (2.8) & 93 (1.8) & 9 (2.1) & -4 (2.1) & -2 (2.1) \\
Llama 4 Maverick & 16 (4.3) & 10 (2.0) & 23 (2.9) & 1 (2.2) & -6 (2.2) & -2 (2.2) \\
Sabi\'a 4 & 31 (5.3) & 46 (3.8) & 77 (3.3) & 5 (2.9) & -7 (2.9) & -1 (3.1) \\
ChatGPT (web) & 1 (1.9) & 18 (2.4) & 16 (2.2) & -13 (1.8) & -5 (1.8) & -1 (2.1) \\
\bottomrule
\end{tabular}

\end{table}
\begin{figure}[htbp]\centering
\includegraphics[width=\textwidth]{17_regression_coefficients.pdf}
\caption{The same coefficients with 95\% intervals, by configuration and term.}
\end{figure}

\clearpage
\section{Ten real exchanges}
Selected deterministically as the shortest qualifying answer for each behaviour, specific-candidate wording, labelled responses. Each block gives the source and record identifier, the exact prompt body and question, the answer verbatim (cut at 1,500 characters), and the codes assigned. Markdown emphasis is stripped and links reduced to their domain in brackets; nothing else is edited.
\input{../../output/reviewed/llm_gpt-5-mini-2025-08-07/tables/examples.tex}

\end{document}
